\begin{figure}[!h]
    \centering
    \begin{overpic}[width=0.75\linewidth]{chapters/ch-experiments/img/seeds.pdf}
    \put(2.5,7){\footnotesize 20}
    \put(2.5,19){\footnotesize 40}
    \put(2.5,32){\footnotesize 60}
    \put(2.5,45){\footnotesize 80}
    \put(10.5,3){\footnotesize 1}
    \put(20,3){\footnotesize 2}
    \put(29,3){\footnotesize 3}
    \put(38,3){\footnotesize 4}
    \put(47.5,3){\footnotesize 5}
    \put(57,3){\footnotesize 6}
    \put(66,3){\footnotesize 7}
    \put(75,3){\footnotesize 8}
    \put(84,3){\footnotesize 9}
    \put(93,3){\footnotesize 10}
    \put(43,-1.8){\footnotesize Incremental Step}
    \put(-2,24){\footnotesize \rotatebox{90}{Acc}}
    \put(88,47.5){\footnotesize 42}
    \put(88,43){\footnotesize 123}
    \put(88,38.7){\footnotesize 1000}
    \put(88,34.5){\footnotesize 1993}
    \put(88,30.5){\footnotesize 2023}
    \end{overpic}
    \caption{Performance following each incremental step with the varied task sequence.% The mean and standard deviation across four splits are represented by solid lines and shaded areas, respectively. Each unique task sequence, determined by a random seed, is distinguished by a different color.
    }
    \label{fig:order_b10}
\end{figure}